\section{Workstream 1: data collection}
\label{sec:data}

Report criterion 1 is worth 5 points, the most of any report line, and it covers both collection
and processing. Every step below therefore records what the report needs: the source, the
method, the date, the numbers, and the problems hit along the way.

\subsection{DS1: the core snapshot}
\label{sec:ds1}

\textbf{Purpose.} The 25 + 3 addresses that the Explorer shows, and the data every method is
evaluated on. This is A1 Objective 1.

\textbf{Selection rule.} It is written into the manifest \emph{before} any history is fetched,
which avoids A1 risk 19, choosing addresses because they already look odd. All parameters come
from \code{settings.json}:
\begin{enumerate}
  \item Pick 20 block heights spread evenly between tip − 4,000 and tip − 100. That covers about
    four weeks of blocks and avoids reorganisations. The seed is 30049.
  \item From each block, take the addresses on the outputs of up to 4 pages of transactions (25
    per page) to form the candidate pool. Duplicates are removed and script types recorded.
  \item Shuffle the pool with seed 30049 and visit candidates in that order. For each, call
    \code{/address/<a>} once (cheap) and read \code{chain_stats.tx_count}.
  \item \textbf{Main addresses:} keep a candidate if it has 50--1,000 confirmed transactions and
    its full history spans at least 90 days (checked after fetching). Stop at 25.
  \item \textbf{Short-history addresses:} keep the first 3 candidates with 3--9 transactions.
    These exercise the insufficient-history state.
  \item Every visited candidate is logged with the reason it was kept or skipped (too few, too
    many, span too short, fetch error). The skip counts go into the report as evidence that
    the rule was applied.
\end{enumerate}

\textbf{Fetching.} \code{python -m explorer.cli snapshot} pages each history with
\fp{/address/<a>/txs/chain[/<last_txid>]}, newest first, 25 per page, until a page holds fewer
than 25 transactions or the cap of 1,000 is reached. It waits at least 0.5 s between requests and
backs off exponentially on 429 and 5xx, honouring \code{Retry-After}. Each raw page is saved
unchanged as \code{data/raw/esplora/<address>/page-NNN.json.gz}, so a re-run reads the cache
instead of the network unless \code{--refetch} is given.

\textbf{Verification}, run automatically after the fetch and printed into the manifest:
\begin{itemize}
  \item Number fetched equals \code{chain_stats.tx_count}.
  \item No duplicate txids, every transaction confirmed, and every block time present.
  \item Summed \code{received} equals \code{funded_txo_sum}, and summed \code{sent} equals
    \code{spent_txo_sum}, for every address (risk 18).
  \item For 2 addresses, the first page from the fallback (mempool.space) matches in time and
    amount.
\end{itemize}

\begin{team}
A1 promised that a second member would verify one transaction's amount and direction by hand
against the raw inputs and outputs, reproduce the retrieval, and review the manifest. Ask Cam
Dan Bui (WBS 1.2.1 co-owner) to do this on 22 October and record their name and the date in the
manifest.
\end{team}

\textbf{Manifest} (\code{data/processed/esplora/manifest.json}), with every DR-01 field: source
and API base, retrieval timestamps, the request pace and the request statistics (count, 429s,
retries, median latency), the tip height, the selection rule and seed, kept and skipped
candidates with reasons, the fields used, missing-value counts, provenance (raw page paths and
SHA-256), the terms note, and rebuild instructions.

\textbf{Expected output}\est{}. 28 addresses and about 5,000 transactions (median about 150 per
address), with spans of 3--24 months.

\subsection{DS2: the sequence corpus for TxGPT}
\label{sec:ds2}

\textbf{Purpose.} A language model needs far more sequences than 25 addresses provide. DS2 is
our own large crawl, and it counts as an additional dataset.

\textbf{Discovery.} The same block-sampling method as DS1, over 200 blocks spread across the
last 12 months, which gives seasonal variety. That yields a pool of around 60,000 distinct
addresses\est{}. DS1's 28 addresses are removed from the pool before anything else happens, so
TxGPT never trains on the addresses it later scores (Section~\ref{sec:leakage}).

\textbf{Inclusion.} 20--1,000 confirmed transactions. The target is 5,000 addresses, the minimum
3,000, the stretch 10,000.

\textbf{Compact storage.} DS2 does not keep raw pages. Each page is reduced at once to one row
per transaction, which keeps DS2 near 60 MB instead of tens of gigabytes:
\begin{center}
\small
\begin{tabularx}{\textwidth}{@{}L{3.3cm}L{2.0cm}Y@{}}
\toprule
\textbf{Column} & \textbf{Type} & \textbf{Meaning} \\
\midrule
\code{addr_id} & string & SHA-256 of the address with a project salt, first 16 hex characters. The salt is kept out of the zip (A1 risk 16). \\
\code{txid_hash} & string & Hash of the txid, used only to remove duplicates. \\
\code{block_height}, \code{block_time_utc} & int & Confirmation height and Unix time. \\
\code{received_sats}, \code{sent_sats} & int & As \code{amounts.py} computes them. \\
\code{n_inputs}, \code{n_outputs} & int & Transaction shape. \\
\code{fee_sats}, \code{vsize_vb} & int & For the fee-rate token. \\
\bottomrule
\end{tabularx}
\end{center}

\textbf{Operations.} A state file records each address as done, partial (with its last txid) or
failed, so the crawl resumes after any interruption. The pace is 0.5 s with back-off. If
Blockstream returns repeated 429s, the client switches to mempool.space for an hour and then
returns. Expected load\est{}: about 5,000 addresses × 4 pages on average gives 20,000 requests at
about 1.2 s each, roughly 7 hours, run overnight on 22 October.

\textbf{Checks.} Every address's row count matches its \code{tx_count}, up to the cap. There are
no duplicate \code{txid_hash} values within an address, and no DS1 address appears.

\subsection{DS3: Elliptic}
\label{sec:ds3}

\begin{itemize}
  \item \textbf{Source.} Kaggle, ``Elliptic Data Set'', released with Weber et al.\ (2019). Three
    CSVs: features (txId, time step, 165 features), classes (\code{1} = illicit, \code{2} =
    licit, \code{unknown}) and the edge list.
  \item \textbf{Download.} \code{python -m explorer.cli external --elliptic}, which uses the
    Kaggle CLI if a token is set and otherwise prints the manual steps. It records SHA-256 of
    each file.
  \item \textbf{Licence.} Read from the Kaggle page on the day of download and copied into the
    data card. If redistribution isn't allowed, the zip carries the download script and
    checksums instead of the data, and the README says so.
  \item \textbf{Challenges for the report.} The features are anonymised, so they can't be
    computed for our own addresses. That is why DS3 is used to evaluate method families, not to
    score DS1. About 77\% of nodes are unlabelled. The authors report that a dark-market shutdown
    around time step 43 changes the data distribution.
\end{itemize}

\subsection{DS4: BitcoinHeist}
\label{sec:ds4}

\begin{itemize}
  \item \textbf{Source.} UCI Machine Learning Repository, dataset 526, from Akcora et al.\
    (2020). A1 already cites the UCI page. Licence CC BY 4.0 (checked 20 October 2026); a
    110.7 MB zip holding one 225 MB CSV: \code{address, year, day, length, weight, count, looped,
    neighbors, income, label}.
  \item \textbf{Download.} \code{python -m explorer.cli external --bitcoinheist}, with SHA-256
    recorded.
  \item \textbf{Challenges for the report.} Severe imbalance (positives are about 1.4\% of
    rows\est{}). Labels come from three earlier ransomware studies, so ``white'' isn't verified
    innocence. Features are 24-hour graph statistics, not single transactions, which is why A1
    called it a comparison source. Family coverage differs by year, which the time-split design
    must respect.
\end{itemize}

\subsection{DS5 and DS6: text corpora}
\label{sec:ds56}

Collected by the existing scripts. \code{python -m lm.prepare_text --delete-raw} streams
FineWeb-Edu shards, trains our 32,768-token BPE tokenizer on the first 100,000 documents, writes
2B training tokens plus 10M validation tokens, and writes \code{meta.json}.
\code{python -m lm.prepare_chat} renders smol-smoltalk with our chat template. Conversations that
are too long are truncated at an assistant turn rather than dropped (about 46\% exceed 1,024
tokens). It writes the token and loss-mask files. Both scripts already record source, revision,
counts and settings, which the data card reuses.

\subsection{DS7: the explanation set}
\label{sec:ds7}

Built on 29 October, after the models exist, by \code{python -m explorer.cli build-ds7}:
\begin{itemize}
  \item \textbf{About 10,000 explanation pairs.} The input is a JSON summary of a scored
    transaction from DS1 or DS2: direction, amount, the usual amount and gap, each feature's z,
    the comparison basis, and which features contribute. The output is the Explorer's template
    explanation, the same wording rule as the A1 detail panel, ending with the fixed limitation
    line. Five phrasings of each template stop the model memorising one sentence.
  \item \textbf{About 300 glossary pairs written by the team}, for example ``What does a robust z
    of 4 mean?'' or ``Why isn't an unusual transaction proof of wrongdoing?''. These are written
    from our own report text.
  \item \textbf{Split by source address}: 90\% train, 5\% validation, 5\% test, so test pairs
    never come from a training address.
\end{itemize}

\subsection{Collection problems and responses, for the report}
\label{sec:challenges}

\begin{center}
\small
\begin{tabularx}{\textwidth}{@{}L{4.2cm}Y@{}}
\toprule
\textbf{Problem} & \textbf{Response} \\
\midrule
Esplora pages hold 25 transactions and rate limits aren't published & Cursor paging; a pause of at least 0.5 s; exponential back-off honouring \code{Retry-After}; raw-page caching; request statistics logged in the manifest. \\
Raw transaction JSON is large (about 1.3 MB per page for busy addresses) & DS1 keeps raw pages gzipped as evidence; DS2 reduces each page to compact rows immediately. \\
Very busy addresses (exchanges) exceed the 1,000-transaction cap & Excluded by the rule; stated as a bias (the sample leans towards ordinary wallets). \\
Address-level amount and direction aren't API fields & Derived from \code{vout} and \code{vin.prevout}; verified by the sum checks and a manual check. \\
Selection bias & Rule fixed before fetching; seeded; all skips logged. \\
No ground-truth labels for Bitcoin addresses & Two labelled public datasets (DS3, DS4) for label-based metrics; planted transactions for known positives. \\
Labelled datasets are anonymised or aggregated & Used to evaluate method families, not to score DS1. \\
Licences and privacy & Licences recorded per dataset; DS2 addresses hashed; no identity claims. \\
\bottomrule
\end{tabularx}
\end{center}

\begin{done}
DS1 holds 25 + 3 addresses and passes every automated check, with the manifest reviewed by a
teammate. DS2 holds at least 3,000 addresses and passes its checks. DS3 and DS4 are downloaded
with checksums and licences recorded. DS5 and DS6 are tokenised. Collection numbers (requests,
429s, skips) are copied into \fp{reports/tables/collection_summary.csv}, ready for the report.
\end{done}
